
où les
flèches verticales sont les homomorphismes canoniques @@M0@@ (IV 2.5.1).

\subsection*{7.7. Spécialisation pour @@M1@@ et @@M2@@.}
Lorsque @@M3@@ et @@M4@@ sont réguliers (par exemple si @@M5@@ est lisse), les deux flèches verticales précédentes sont des isomorphismes (IV 2.5), et l'homomorphisme de spécialisation peut alors s'interpréter comme un homomorphisme
@@M6@@
qui peut se décrire encore comme l'unique homomorphisme rendant commutatif le diagramme suivant
@@M7@@
Comme dans ce dernier @@M8@@ et @@M9@@ sont des homomorphismes d'anneaux (IV 2.7 b) avec @@M10@@ surjectif, il s'ensuit que \emph{l'homomorphisme de spécialisation} (7.7.1) \emph{est un homomorphisme d'anneaux}. Pour la même raison, \emph{si} @@M11@@ \emph{est séparé}, \emph{l'homomorphisme} (7.7.1) \emph{est même un homomorphisme de} @@M12@@-\emph{anneaux}. Pour préciser ce point, il faut seulement noter qu'un schéma noethérien régulier séparé est divisoriel (II 2.2.7.1), donc son @@M13@@ s'identifie à son @@M14@@ «~naïf~» (IV 2.9 (ii)), lequel est muni de façon naturelle d'une structure de @@M15@@-anneau (VI 3.2). D'autre part, le fait que @@M16@@ et @@M17@@ soient réguliers implique le même fait pour @@M18@@ (EGA @@M19@@ 17.1.8), et les homomorphismes @@M20@@ et @@M21@@ sont des @@M22@@-homomorphismes, donc il en est de même de (7.7.1) qui se déduit de @@M23@@ par passage au quotient.

